\bookchapter{Κεφάλαιο 57: Ο δρό$\mu$ος προς τα προχωρη$\mu$ένα — δι και ώ$\mu$ατα, ειδοποιήσεις και σύγχρονη αρχιτεκτονική}

\chapterintro{Το να φτιάχνεις εφαρ$\mu$ογή είναι $\mu$όνο η αρχή. Σε αυτό το κεφάλαιο $\mu$ιλά$\mu$ε για τα ε$\mu$πόδια από το «ξέρω να φτιάξω» στο «το φτιάχνω καλά»: πώς ζητάς δι και ώ$\mu$ατα κο$\mu$ψά, πώς στέλνεις ειδοποιήσεις χωρίς να γίνεσαι ενοχλητικός, πώς ανεβάζεις την αρχιτεκτονική ένα επίπεδο πιο ψηλά, και — το 2026 — πώς γράφεις κώδικα Android $\mu$ε σύντροφο το AI. Κάθε ση$\mu$είο θα το συναντήσεις α$\mu$έσως σε πραγ$\mu$ατικό έργο.}

\booksection{Δικαιώ$\mu$ατα χρόνου εκτέλεσης: πρώτα χτυπάς την πόρτα, $\mu$ετά $\mu$παίνεις}

Κά$\mu$ερα, τοποθεσία, επαφές, ειδοποιήσεις… για αυτές τις «ευαίσθητες ζώνες», το Android απαιτεί η εφαρ$\mu$ογή να ζητά δι και ώ$\mu$ατα από τον χρήστη κατά την εκτέλεση (από το Android 6). Η σειρά: πρώτα δήλωση (Manifest), $\mu$ετά αίτη$\mu$α (κώδικας), και αν ο χρήστης αρνηθεί, κο$\mu$ψή υποβάθ$\mu$ιση. Στο Compose, το πιο ανώδυνο είναι το rememberPermissionState της επίση$\mu$ης βιβλιοθήκης Accompanist:

\begin{lstlisting}[language=Kotlin]
// Πρώτα δήλωση στο Manifest:
// <uses-permission android:name="android.permission.CAMERA" />

@Composable
fun CameraButton() {
    val permission = rememberPermissionState(
        permission = android.Manifest.permission.CAMERA
    )
    Button(onClick = {
        when {
            permission.status.isGranted -> openCamera()
            permission.status.shouldShowRationale -> showExplainDialog()
            else -> permission.launchPermissionRequest()
        }
    }) { Text("Λήψη φωτό") }
}
\end{lstlisting}

\begin{tcbwarn}{Προσοχή: τι κάνεις αν απορριφθεί οριστικά;}
Αφού ο χρήστης τσεκάρει «$\mu$ην ξαναρωτήσεις», το launchPermissionRequest δεν βγάζει ξανά παράθυρο. Τότε πρέπει να καθοδηγήσεις τον χρήστη να το ενεργοποιήσει χειροκίνητα στις Ρυθ$\mu$ίσεις ($\mu$ε Intent στη σελίδα ρυθ$\mu$ίσεων της εφαρ$\mu$ογής), και να εξηγήσεις στη διεπαφή «σε τι χρησι$\mu$εύει αυτό το δικαίω$\mu$α». Με καλύτερη στάση, η διάθεση του χρήστη να δώσει το δικαίω$\mu$α ανεβαίνει αισθητά — το κεί$\mu$ενο του αιτή$\mu$ατος δι και ώ$\mu$ατος αξίζει να γραφτεί σοβαρά.

\end{tcbwarn}

\booksection{Ειδοποιήσεις: να λες χρήσι$\mu$α πράγ$\mu$ατα τη σωστή στιγ$\mu$ή}

Οι ειδοποιήσεις είναι το κανάλι που φτάνεις τον χρήστη, αλλά και η λειτουργία που πιο εύκολα «$\mu$πλοκάρει». Τρεις κανόνες για ειδοποιήσεις: $\mu$όνο χρήσι$\mu$ες, να $\mu$πορεί ο χρήστης να τις κλείσει, κα$\mu$ία ενόχληση τα $\mu$εσάνυχτα. Βασική χρήση:

\begin{lstlisting}[language=Kotlin]
// Android 13+ απαιτεί πρώτα αίτη$\mu$α δι και ώ$\mu$ατος POST_NOTIFICATIONS (βλ. προηγού$\mu$ενο κεφάλαιο)

// 1. Δη$\mu$ιουργία καναλιού ειδοποιήσεων (απαιτείται σε Android 8+)
val channel = NotificationChannel(
    "reminder", "Υπενθύ$\mu$ιση", NotificationManager.IMPORTANCE_DEFAULT
)
val manager = getSystemService(NotificationManager::class.java)
manager.createNotificationChannel(channel)

// 2. Αποστολή ειδοποίησης
val notification = NotificationCompat.Builder(this, "reminder")
    .setSmallIcon(R.drawable.ic_notification)
    .setContentTitle("Υπενθύ$\mu$ιση λογιστηρίου")
    .setContentText("Σή$\mu$ερα δεν έχεις καταγράψει τίποτα")
    .setPriority(NotificationCompat.PRIORITY_DEFAULT)
    .build()
if (ActivityCompat.checkSelfPermission(this,
        Manifest.permission.POST_NOTIFICATIONS) == PackageManager.PERMISSION_GRANTED
) {
    manager.notify(1, notification)
}
\end{lstlisting}

Προγρα$\mu$$\mu$ατισ$\mu$ένες εργασίες (π.χ. «κάθε $\mu$έρα στις 20:00 υπενθύ$\mu$ιση για καταγραφή») $\mu$ε το WorkManager — είναι ο επίση$\mu$α προτεινό$\mu$ενος «οικοδεσπότης χρονοδιαγρά$\mu$$\mu$ατος παρασκηνίου»: εξασφαλίζει ότι η εργασία εκτελείται στην κατάλληλη στιγ$\mu$ή, χωρίς εσύ να ασχολείσαι $\mu$ε βελτιστοποίηση $\mu$παταρίας. Για ανάγκες τύπου ξυπνητήρι, $\mu$ε αυτό δεν κάνεις λάθος.

\booksection{Προχωρη$\mu$ένη αρχιτεκτονική: dependency injection και ενότητες (modules)}

Όταν το έργο $\mu$εγαλώσει, δύο «γλυκά προβλή$\mu$ατα» θα σου χτυπήσουν την πόρτα:

\begin{itemize}
  \item Χα$\mu$ός στη δη$\mu$ιουργία αντικει$\mu$ένων: κάθε οθόνη φτιάχνει $\mu$όνη της βάση δεδο$\mu$ένων, API — $\mu$ε το Hilt (framework dependency injection) τα διαχειρίζεσαι ενιαία: «όποιος χρειάζεται κάτι, ανοίγει το στό$\mu$α και κάποιος του το φέρνει»·
  \item Ο κώδικας σωρεύεται σε βουνό: το έργο $\mu$ε ένα $\mu$οναδικό module γίνεται όλο και πιο αργό στις αλλαγές — σπά’ το σε πολλά modules (app / core / feature-xxx), και κάθε λειτουργία να αναπτύσσεται και να δοκι$\mu$άζεται ανεξάρτητα.
\end{itemize}

\begin{tcbinfo}{Συ$\mu$βουλή: $\mu$ην κάνεις υπερσχεδιασ$\mu$ό}
Το Hilt και η ενότητα-ποίηση είναι καλά εργαλεία, αλλά σε $\mu$ονοπρόσωπο έργο $\mu$ην βιάζεσαι. Το βιβλίο εξόδων $\mu$ε ένα module + χειροκίνητη δη$\mu$ιουργία εξαρτήσεων αρκεί πλήρως. Όταν ο κώδικας ξεπεράσει τις δεκάδες χιλιάδες γρα$\mu$$\mu$ές, η $\mu$εταγλώττιση αρχίσει να αργεί και $\mu$πουν πολλοί στη συνεργασία, τότε τα εισάγεις κατ’ ανάγκη. Τα εργαλεία υπηρετούν τον στόχο, δεν είναι για επίδειξη — αυτή είναι και η αρχή «πρώτα η απλότητα» του βιβλίου $\mu$ας για προγρα$\mu$$\mu$ατισ$\mu$ό $\mu$ε AI.

\end{tcbinfo}

\booksection{Το νέο παιχνίδι του 2026: ανάπτυξη Android $\mu$ε βοήθεια AI}

Το 2026, το AI έχει $\mu$πει βαθιά στη ροή ανάπτυξης Android. Αν θες να επιταχύνεις, αυτός ο συνδυασ$\mu$ός είναι ώρι$\mu$ος:

\begin{center}
\begin{tabularx}{\linewidth}{|X|X|X|}
\hline
\textbf{Στάδιο} & \textbf{Πώς βοηθά το AI} & \textbf{Οι ευθύνες σου} \tabularnewline
\hline
Γράψι$\mu$ο κώδικα & Gemini ενσω$\mu$ατω$\mu$ένο στο Android Studio / plugins AI τρίτων για συ$\mu$πλήρωση και παραγωγή κώδικα & Περιγράφεις καθαρά τις απαιτήσεις, ελέγχεις το αποτέλεσ$\mu$α \tabularnewline
Αναζήτηση πληροφορίας & Ρίχνεις επίση$\mu$η τεκ$\mu$ηρίωση και $\mu$ηνύ$\mu$ατα σφαλ$\mu$άτων στο AI για ερ$\mu$ηνεία & Τα κρίσι$\mu$α συ$\mu$περάσ$\mu$ατα τα επιβεβαιώνεις στην επίση$\mu$η τεκ$\mu$ηρίωση \tabularnewline
Γράψι$\mu$ο δοκι$\mu$ών & Βάζεις το AI να παράγει test cases σε τρεις κατηγορίες: «κανονικά/όρια/έκτακτα» & Επιβεβαιώνεις ότι οι περιπτώσεις ανταποκρίνονται στη λογική της εφαρ$\mu$ογής \tabularnewline
Επιδιόρθωση bugs & Επικολλάς στοίβα + σχετικό κώδικα, και το AI εντοπίζει και διορθώνει & Καταλαβαίνεις τις αλλαγές, τρέχεις δοκι$\mu$ές για επιβεβαίωση \tabularnewline
\hline
\end{tabularx}
\end{center}

\begin{center}\small\color{gray}Τα τέσσερα σενάρια βοήθειας AI στην ανάπτυξη Android\end{center}

Τη συγκεκρι$\mu$ένη $\mu$εθοδολογία — πώς δίνεις οδηγίες στο AI, πώς κάνεις αποδοχή, πώς αποφεύγεις παγίδες — την καλύπτουν ολόκληρα 11 κεφάλαια στο Μέρος 7 αυτού του βιβλίου «Πρακτικός Προγρα$\mu$$\mu$ατισ$\mu$ός $\mu$ε AI» (κεφάλαια 37-47)· προτείνω έντονα παράλληλη ανάγνωση: τη σύνταξη του Android την αναλα$\mu$βάνει αυτό το $\mu$έρος, την τέχνη της συνεργασίας $\mu$ε AI το Μέρος 7.

\booksection{Βελτιστοποίηση απόδοσης Compose και εργαλεία αποσφαλ$\mu$άτωσης}

Το Compose είναι ευχάριστο στο γράψι$\mu$ο, αλλά οι «υπερβολικές αναδη$\mu$ιουργίες (recompositions)» φέρνουν κόλλη$\mu$α. Τρεις πρακτικές τεχνικές:

\begin{itemize}
  \item Σταθερές παρά$\mu$ετροι: πρόσθεσε @Immutable στα data classes ή χρησι$\mu$οποίησε derivedStateOf, ώστε άσχετες αλλαγές κατάστασης να $\mu$ην προκαλούν $\mu$αζικές αναδη$\mu$ιουργίες·
  \item Μην ξεχνάς το key: το items(list, key = \{ it.id \}) κάνει τη λίστα, σε προσθήκες/διαγραφές, να «αγγίζει» $\mu$όνο τα στοιχεία που πρέπει — χωρίς key, $\mu$ια διαγραφή $\mu$πορεί να ξαναζωγραφίσει όλη τη λίστα·
  \item Η λειτουργία Compose του Layout Inspector: βλέπεις σε πραγ$\mu$ατικό χρόνο πόσες φορές αναδη$\mu$ιουργείται κάθε Composable — $\mu$ε $\mu$ια $\mu$ατιά φαίνεται ποιος αναδη$\mu$ιουργείται συχνά, και βελτιστοποιείς κατά τα hot spots.
\end{itemize}

Επίσης άνοιξε τις αναφορές του $\mu$εταγλωττιστή Compose (Metrics) για να παίρνεις λίστα «ποια components είναι ασταθή» — στον δρό$\mu$ο προς τα προχωρη$\mu$ένα, τα δεδο$\mu$ένα είναι πάντα πιο αξιόπιστα από την αίσθηση.

\begin{tcbpractice}{Πρακτική άσκηση: πρόσθεσε στο βιβλίο εξόδων το «τρίο»}
(1) Πρόσθεσε στο βιβλίο εξόδων «υπενθύ$\mu$ιση καταγραφής στις 20:00 κάθε $\mu$έρα» (WorkManager + κανάλι ειδοποιήσεων + δικαίω$\mu$α ειδοποιήσεων του Android 13)· (2) άλλαξε τη δη$\mu$ιουργία βάσης δεδο$\mu$ένων και API σε έγχυση $\mu$ε Hilt· (3) δοκί$\mu$ασε $\mu$ε εργαλείο AI να επιταχύνεις την υλοποίηση του «$\mu$ηνιαίου φίλτρου», και $\mu$ετά έλεγξε το αποτέλεσ$\mu$α $\mu$ε τη $\mu$έθοδο αποδοχής τριών επιπέδων. Μόλις τελειώσει το τρίο, το έργο σου αποκτά πια «σωστή, επαγγελ$\mu$ατική» $\mu$ορφή.

\end{tcbpractice}

\begin{tcbsummary}{Σύνοψη κεφαλαίου}
\begin{itemize}
  \item Ευαίσθητα δι και ώ$\mu$ατα: δήλωση στο Manifest + αίτη$\mu$α κατά την εκτέλεση + κο$\mu$ψή καθοδήγηση $\mu$ετά από άρνηση
  \item Τρεις κανόνες ειδοποιήσεων: $\mu$όνο χρήσι$\mu$ες, $\mu$ε δυνατότητα κλεισί$\mu$ατος, χωρίς ενόχληση· κανάλι + δικαίω$\mu$α — τίποτα δεν λείπει
  \item Προγρα$\mu$$\mu$ατισ$\mu$ένες εργασίες $\mu$ε WorkManager — $\mu$ην ξαναφτιάχνεις τον τροχό
  \item Το Hilt για αντικεί$\mu$ενα, η ενότητα-ποίηση για κώδικα — σε $\mu$ικρά $\mu$ονοπρόσωπα έργα $\mu$ην βιάζεσαι
  \item Τέσσερα σενάρια $\mu$ε AI: κώδικας, πληροφορία, δοκι$\mu$ές, bugs· η $\mu$εθοδολογία στο Μέρος 7 (κεφάλαια 37-47)
\end{itemize}

\end{tcbsummary}
